% Chapter 2: Study and design. Skeleton: each \plan{} says what the section will hold and
% where its facts come from. Prose replaces the plans section by section.

\chapter{Study and design}
\label{ch:design}

\section{Introduction}
\plan{Half a page. What this chapter does: turn the problem from the general introduction
into a specification, show the system architecture, then present every electronic component
with its working principle and the reason it was chosen. The modular platform comes last.}

\section{Specification}
\label{sec:spec}

\subsection{Origin of the need}
BLOC-BOT grew out of our own experience of student robotics competitions. Our team took
part in two consecutive editions of the Algerian Robot Cup (ARC~3 and ARC~4), and in
several other competitions, without success. Each ARC edition has two stages. To
qualify, a team's robot must follow a line through a maze, and the fastest robots go
through. The main event then calls for a second robot, built for a task that changes
from one year to the next. Looking back, we traced most of our difficulties to two
causes.

The first was time. For every competition we built the robot's modules from scratch:
a new chassis, new sensor mounts, new wiring, and in the case of ARC, two robots instead
of one. Each took long enough that we reached the competition underprepared, with too
little time left for testing and tuning.

The second was the specification itself. Organisers often announce the rules late or in
vague terms. Unsure of what would be asked, we tended to over-prepare and build a robot
more complicated than the task required. A more complicated robot has more ways to fail,
and some of ours failed during the competition itself.

Both problems point to the same answer: a single robot that is ready in advance and can
be reconfigured quickly for each competition, rather than a new robot every time. Most
competitions in this category rely on a small, recurring set of sensors (line sensors,
distance sensors, wheel encoders). If each of these has a standard way of being mounted
and wired, adapting the robot becomes a matter of choosing modules instead of designing
them. This is the idea behind BLOC-BOT and its interchangeable mechanical holders, which
we call \emph{blocs}.

\subsection{Target users}
The platform is aimed first at students and teams preparing for robotics competitions,
who need a working base quickly and want to spend their time on strategy and tuning rather
than on rebuilding mechanics. Because the hardware and firmware are open and the
platform is designed to be extended, it is also meant for students and researchers who
need a small mobile robot as a test bench for their own sensors or algorithms.

\subsection{Functional requirements}
Table~\ref{tab:requirements} lists what the system must do. The last column gives the
state of each requirement at the time of writing, so that the specification can be read
against the results in Chapter~\ref{ch:realisation}.

\begin{table}[htbp]
  \centering
  \caption{Functional requirements of BLOC-BOT and their state at the time of writing.}
  \label{tab:requirements}
  \begin{tabularx}{\textwidth}{@{}l X l@{}}
    \toprule
    ID & Requirement & State \\
    \midrule
    FR1 & The robot is driven by hand from a handheld controller over a wireless link,
          with no computer or phone in the loop. & Done \\
    FR2 & The controller shows feedback from the robot, starting with the speed of its
          wheels. & Done \\
    FR3 & The robot stops on its own when it no longer receives commands from the
          controller. & Done \\
    FR4 & Sensors and modules are fixed to the robot through a standard hole pattern and
          printed holders, without redesigning the chassis. & Done \\
    FR5 & The controller runs from its own rechargeable battery. & Done \\
    FR6 & The robot follows a line using a multi-channel infrared sensor. & In progress \\
    FR7 & The robot detects obstacles and parks itself using distance sensors. & Planned \\
    \bottomrule
  \end{tabularx}
\end{table}

\subsection{Non-functional requirements and constraints}
Besides what the system does, four conditions shaped how it was built.

\begin{itemize}
  \item \textbf{Local availability of components.} Most of the components were chosen
        because they could be bought in Algeria or were already in our hands, not after a
        comparison of every alternative on the market. We state this plainly here, and
        each component section in this chapter says when availability was the deciding
        reason.
  \item \textbf{Low cost.} The platform is meant for students, so it is built from common
        hobbyist modules and 3D-printed parts. Its cost is detailed in Chapter~3.
  \item \textbf{Reproducibility.} The mechanical parts are printed on an ordinary
        desktop fused deposition modelling (FDM) printer, and the hardware designs and firmware are to be published
        as open source.
  \item \textbf{Size.} The size is set by the waffle plate, a
        \(210 \times 210\)~mm printed plate that carries every other part. In the CAD
        model of the assembled robot, the wheels stand out on each side, giving an
        overall width of 242~mm and a length of 214~mm. The robot is 102~mm tall up to
        its upper plate and 133~mm with the ultrasonic sensor mounted on top (values
        measured in the Fusion assembly, version 13).
\end{itemize}

\section{System architecture}
\label{sec:architecture}
BLOC-BOT is made of two separate boards that share no wire: a handheld controller and the
robot itself. Each has its own microcontroller and its own battery, and the two talk
only over a wireless ESP-NOW link (Section~\ref{sec:espnow}). Figure~\ref{fig:architecture}
shows both, with the main signals between their parts.

The \textbf{controller} is the human side of the system. It reads the joystick and the
four buttons, shows a menu and the robot's state on a small OLED screen, and sends drive
commands to the robot every 50~ms while its Drive screen is open.

The \textbf{robot} turns each drive command into a speed for each wheel and drives the two
TT gearmotors through a DRV8833 dual H-bridge. Two LM393 slotted sensors count the turns
of the wheels, and every 100~ms the robot sends the measured wheel speeds back to the
controller, which displays them. If drive commands stop arriving, the robot stops its
motors on its own, which meets requirement FR3. The robot also carries an HC-SR04
ultrasonic sensor at its front, which measures the distance to obstacles, and an MPU6050
inertial sensor, which measures its tilt and rotation. Both readings are sent to the
controller, the distance every 100~ms and the tilt every 50~ms, and shown on their own
screens.

The two power chains are independent. The robot runs from a 2S lithium-ion pack that
feeds the motor driver directly and the microcontroller through a 5~V step-down
converter. The controller runs from a single lithium-ion cell with its own charger and a
5~V step-up converter. Both are detailed in Sections~\ref{sec:robot}
and~\ref{sec:controller}.

\begin{figure}[htbp]
  \centering
  \begin{tikzpicture}[
      font=\footnotesize,
      >={Stealth[length=2mm]},
      block/.style={draw, rounded corners=2pt, align=center, minimum height=8mm,
                    minimum width=26mm, fill=white},
      mcu/.style={block, fill=gray!15, very thick, minimum width=34mm, minimum height=22mm},
      pwr/.style={block, fill=red!7},
      future/.style={block, dashed},
      lbl/.style={font=\scriptsize, midway, fill=white, inner sep=1pt},
      frame/.style={draw=gray!70, rounded corners=4pt, inner sep=3mm},
    ]
    % ---------------- controller ----------------
    \node[mcu] (mcuC) at (0,0) {\textbf{ESP32 DevKit}\\controller};
    \node[block] (joy) at (-5,0.6) {KY-023 joystick};
    \node[block] (btn) at (-5,-0.6) {4 push buttons};
    \node[block] (oled) at (5,0) {OLED 0.96'' 128$\times$64};
    \node[pwr, minimum width=20mm] (cellC) at (-5.4,2.3) {18650 cell};
    \node[pwr, minimum width=22mm] (tp) at (-2.6,2.3) {TP4056 charger};
    \node[pwr, minimum width=12mm] (sw) at (-0.45,2.3) {switch};
    \node[pwr, minimum width=24mm] (boost) at (2.0,2.3) {MT3608 boost};

    \draw[->] (joy) -- (mcuC.west |- joy) node[lbl] {2 $\times$ analog};
    \draw[->] (btn) -- (mcuC.west |- btn) node[lbl] {digital};
    \draw[->] (mcuC.east |- oled) -- (oled) node[lbl] {I$^2$C};
    \draw[->, red!60!black, thick] (cellC) -- (tp);
    \draw[->, red!60!black, thick] (tp) -- (sw);
    \draw[->, red!60!black, thick] (sw) -- (boost);
    \draw[->, red!60!black, thick] (boost.south) |- ++(0,-0.35) -| ([xshift=6mm]mcuC.north)
      node[pos=0.75, right, font=\scriptsize] {5 V};

    \begin{scope}[on background layer]
      \node[frame, fit=(joy)(btn)(oled)(cellC)(boost)(mcuC),
            label={[font=\small\bfseries, anchor=south west]north west:Controller}] (fC) {};
    \end{scope}

    % ---------------- robot ----------------
    \node[mcu] (mcuR) at (0,-5.8) {\textbf{ESP32-S3}\\robot};
    \node[block, minimum height=7mm] (enc) at (-4.6,-4.9) {2 $\times$ LM393 slotted};
    \node[block, minimum height=7mm] (us) at (-4.6,-5.85) {HC-SR04};
    \node[block, minimum height=7mm] (mpu) at (-4.6,-6.8) {MPU6050 IMU};
    \node[future] (line) at (-4.6,-8.0) {TCRT5000 line\\sensor (in progress)};
    \node[block] (drv) at (4.5,-5.8) {DRV8833\\dual H-bridge};
    \node[block, minimum width=17mm] (mL) at (7.4,-5.2) {TT motor L};
    \node[block, minimum width=17mm] (mR) at (7.4,-6.4) {TT motor R};
    \node[pwr] (buck) at (1.5,-8.1) {Mini360 buck};
    \node[pwr] (batt) at (4.5,-8.1) {2S 18650 pack};

    \draw[->] (enc) -- (mcuR.west |- enc) node[lbl] {pulses};
    \draw[<->] (mcuR.west |- us) -- (us) node[lbl] {trig / echo};
    \draw[<->] (mcuR.west |- mpu) -- (mpu) node[lbl] {I$^2$C};
    \draw[->, dashed] (line.east) -| ([xshift=-10mm]mcuR.south);
    \draw[->] (mcuR) -- (drv) node[midway, above, font=\scriptsize] {4 $\times$ PWM};
    \draw[->] (drv.east) -- (mL.west);
    \draw[->] (drv.east) -- (mR.west);
    \draw[->, red!60!black, thick] (batt) -- (buck);
    \draw[->, red!60!black, thick] (buck.north) -- (buck.north |- mcuR.south) node[lbl] {5 V};
    \draw[->, red!60!black, thick] (batt) -- (drv) node[lbl] {VM};

    \begin{scope}[on background layer]
      \node[frame, fit=(enc)(line)(mpu)(mR)(mL)(batt)(mcuR),
            label={[font=\small\bfseries, anchor=south west]north west:Robot}] (fR) {};
    \end{scope}

    % ---------------- radio link ----------------
    \draw[->, very thick, blue!60!black] ([xshift=-5mm]mcuC.south) -- ([xshift=-5mm]mcuR.north)
      node[midway, left, align=right, font=\scriptsize] {drive commands\\every 50 ms};
    \draw[<-, very thick, blue!60!black] ([xshift=5mm]mcuC.south) -- ([xshift=5mm]mcuR.north)
      node[midway, right, align=left, font=\scriptsize] {wheel speeds, distance:\\every 100 ms\\
        tilt (IMU): every 50 ms};
    \node[font=\small\bfseries, blue!60!black, align=center]
      at ($(mcuC.south)!0.5!(mcuR.north) + (-4.4,0)$) {ESP-NOW\\wireless link};
  \end{tikzpicture}
  \caption{System architecture. The controller (top) and the robot (bottom) share no
           wire; over ESP-NOW the controller sends drive commands and the robot sends back
           its wheel speeds, the measured distance and its tilt. Red arrows
           are power, black arrows are signals, and the dashed block is in progress.}
  \label{fig:architecture}
\end{figure}

\section{The robot}
\label{sec:robot}

\subsection{Microcontroller: ESP32-S3 (N16R8)}
The robot is controlled by an ESP32-S3, a system on a chip (SoC) from Espressif that
combines a microcontroller with a 2.4~GHz radio. Its main features are summarised in
Table~\ref{tab:esp32s3}. The chip has a dual-core Xtensa LX7 processor running at up to
240~MHz, 512~KB of internal SRAM, and both Wi-Fi (IEEE 802.11b/g/n) and Bluetooth~5
Low Energy \cite{esp32s3_datasheet}. The radio is what makes the wireless link of
Section~\ref{sec:espnow} possible without any extra module.

The chip is not used on its own but inside a module marked ESP-32-S3-N16R8. This is a
third-party module rather than an Espressif one, but it follows the pinout and the
naming of Espressif's ESP32-S3-WROOM-1 \cite{esp32s3_wroom1_datasheet}: N16R8 stands for
16~MB of flash memory and 8~MB of octal pseudo-static RAM (PSRAM). The module adds these
memories, the crystal and a printed antenna, and exposes 36 general-purpose input/output
(GPIO) pins. It sits on a development board laid out like Espressif's ESP32-S3-DevKitC-1,
which brings the pins out to two headers, carries two USB-C ports for programming and
power, and a regulator that turns a 5~V supply into the 3.3~V the module needs
\cite{esp32s3_devkitc1_guide}. On the robot, that 5~V comes from the
step-down converter described in Section~\ref{sec:robot-power}.

\begin{table}[htbp]
  \centering
  \caption{Main features of the ESP32-S3 and of the N16R8 module used on the robot,
           whose figures follow Espressif's WROOM-1 N16R8
           \cite{esp32s3_datasheet,esp32s3_wroom1_datasheet}.}
  \label{tab:esp32s3}
  \begin{tabularx}{\textwidth}{@{}l X@{}}
    \toprule
    Feature & Value \\
    \midrule
    Processor & Dual-core Xtensa 32-bit LX7, up to 240~MHz \\
    Memory & 512~KB SRAM, 384~KB ROM (chip); 16~MB flash, 8~MB PSRAM (module) \\
    Radio & 2.4~GHz Wi-Fi (IEEE 802.11b/g/n), Bluetooth~5 (LE) \\
    GPIOs & 36 on the module, of which GPIO35--37 are taken by the PSRAM \\
    Analog inputs & Two 12-bit SAR ADCs, up to 20 channels; ADC2 cannot be used
                    while Wi-Fi is active \\
    PWM & LED PWM controller, up to 8 channels \\
    Supply & 3.0--3.6~V (module); 5~V on the development board, regulated to 3.3~V \\
    \bottomrule
  \end{tabularx}
\end{table}

\paragraph{Why the ESP32-S3.} The robot is meant to grow: every new bloc brings a sensor,
and every sensor needs pins. The ESP32-S3 was chosen for its number of usable GPIOs, which
leaves room for the line sensor, the ultrasonic sensor and later additions without
changing the microcontroller. Its PSRAM is not used by the current firmware.

\paragraph{Pins that cannot be used freely.} Not all 36 GPIOs are available. With an
octal PSRAM, GPIO35, GPIO36 and GPIO37 carry the memory's data lines and cannot be used
for anything else \cite{esp32s3_datasheet,esp32s3_wroom1_datasheet,esp32s3_devkitc1_guide},
even though our board still brings them out to its header.
GPIO0, GPIO3, GPIO45 and GPIO46 are strapping pins, read by the chip at reset to choose
its boot mode, so nothing connected to them should pull them high or low at power-up.
GPIO19 and GPIO20 carry the USB interface and GPIO43 and GPIO44 the serial port used for
programming \cite{esp32s3_datasheet}. Finally, because the robot uses its radio all the
time, analog sensors must be wired to ADC1 pins (GPIO1 to GPIO10): the ADC2 channels
cannot be read while Wi-Fi is running \cite{esp32s3_datasheet}. This matters for the
five-channel line sensor, which gives analog outputs.

Table~\ref{tab:robot-pins} gives the pins the robot uses at the time of writing. The
DRV8833's sleep input is tied to 3.3~V, so the driver is always enabled and needs no pin.

\begin{table}[htbp]
  \centering
  \caption{GPIO assignment on the robot (firmware file \texttt{pins\_esp32s3.h}).}
  \label{tab:robot-pins}
  \begin{tabularx}{\textwidth}{@{}l l X@{}}
    \toprule
    GPIO & Connected to & Note \\
    \midrule
    13, 16 & DRV8833 AIN1, AIN2 & left motor, PWM \\
    18, 17 & DRV8833 BIN1, BIN2 & right motor, PWM; the pair is swapped in firmware so
                                   that both wheels turn forward together \\
    40 & left LM393 sensor & digital pulses, read by interrupt \\
    41 & right LM393 sensor & digital pulses, read by interrupt \\
    11 & HC-SR04 TRIG & 10~\textmu s trigger pulse \\
    12 & HC-SR04 ECHO & through a 1~k\(\Omega\)/2~k\(\Omega\) divider
                         (Section~\ref{sec:hcsr04}), timed by interrupt \\
    8, 9 & I\textsuperscript{2}C SDA, SCL & MPU6050 at 0x68
                         (Section~\ref{sec:mpu6050}); reserved for more
                         I\textsuperscript{2}C sensors \\
    \bottomrule
  \end{tabularx}
\end{table}

\subsection{Motor driver: DRV8833}
\label{sec:drv8833}
A microcontroller pin can supply only a few milliamperes, far less than a motor draws,
and it cannot reverse the polarity of the voltage across the motor. Both jobs are done by
a motor driver built around an \emph{H-bridge}.

\paragraph{The H-bridge.} An H-bridge is made of four switches arranged around the motor,
as in Figure~\ref{fig:hbridge}. Closing \(Q_1\) and \(Q_4\) connects the left terminal of
the motor to the supply and the right terminal to ground, and the motor turns one way.
Closing \(Q_3\) and \(Q_2\) reverses the voltage, and the motor turns the other way.
Closing both low-side switches \(Q_2\) and \(Q_4\) shorts the motor terminals together,
which brakes it, and opening all four leaves it free to coast. The two switches of the
same leg (\(Q_1\) and \(Q_2\), or \(Q_3\) and \(Q_4\)) must never be closed together,
since that would short the supply to ground; the driver chip prevents this internally.

\begin{figure}[htbp]
  \centering
  \begin{circuitikz}[scale=0.9, transform shape]
    \draw (0,4) -- (4,4);
    \draw (0,0) -- (4,0);
    \draw (2,4) -- (2,4.5) node[above] {\(V_M\)};
    \draw (2,0) node[ground] {};
    \draw (0,4) to[nos, l_=\(Q_1\)] (0,2) to[nos, l_=\(Q_2\)] (0,0);
    \draw (4,4) to[nos, l=\(Q_3\)] (4,2) to[nos, l=\(Q_4\)] (4,0);
    \node[elmech] (mot) at (2,2) {M};
    \draw (0,2) node[circ] {} node[left=12pt] {OUT1} -- (mot.west);
    \draw (mot.east) -- (4,2) node[circ] {} node[right=12pt] {OUT2};
  \end{circuitikz}
  \caption{Principle of an H-bridge. The DRV8833 contains two of these, built from
           N-channel MOSFETs, one per motor \cite{drv8833_datasheet}.}
  \label{fig:hbridge}
\end{figure}

\paragraph{The DRV8833.} The robot uses one DRV8833 from Texas Instruments, a chip that
contains two complete H-bridges, so a single module drives both wheels. Each bridge is
controlled by two logic inputs, \texttt{IN1} and \texttt{IN2}, following
Table~\ref{tab:drv8833-logic}. Its main ratings are listed in
Table~\ref{tab:drv8833}.

\begin{table}[htbp]
  \centering
  \caption{DRV8833 H-bridge logic, per bridge \cite{drv8833_datasheet}, and the
           switches of Figure~\ref{fig:hbridge} that it closes.}
  \label{tab:drv8833-logic}
  \begin{tabular}{@{}cccccl@{}}
    \toprule
    IN1 & IN2 & OUT1 & OUT2 & Switches closed & Function \\
    \midrule
    0 & 0 & Z & Z & none & Coast \\
    1 & 0 & H & L & \(Q_1\), \(Q_4\) & Forward \\
    0 & 1 & L & H & \(Q_3\), \(Q_2\) & Reverse \\
    1 & 1 & L & L & \(Q_2\), \(Q_4\) & Brake \\
    \bottomrule
  \end{tabular}
\end{table}

\begin{table}[htbp]
  \centering
  \caption{Main ratings of the DRV8833 \cite{drv8833_datasheet}.}
  \label{tab:drv8833}
  \begin{tabularx}{\textwidth}{@{}l X@{}}
    \toprule
    Parameter & Value \\
    \midrule
    Motor supply voltage \(V_M\) & 2.7 to 10.8~V \\
    Output current per bridge (per motor) & 1.5~A RMS, 2~A peak (HTSSOP package with
                                 exposed pad); 0.5~A RMS, 2~A peak (plain TSSOP package) \\
    Undervoltage lockout (voltage) & below 2.6~V (\(V_M\) falling, maximum) \\
    Overcurrent protection (current) & trips between 2 and 3.3~A; the bridge switches off and retries
                             after 1.35~ms \\
    Logic input voltage & high from 2~V (\texttt{nSLEEP} from 2.5~V), up to 5.75~V \\
    \bottomrule
  \end{tabularx}
\end{table}

\paragraph{Speed control by PWM.} To vary the speed, the firmware does not hold an input
high but switches it on and off rapidly, a technique called pulse-width modulation (PWM).
If the input is high for a fraction \(D\) of each period, the duty cycle, the motor sees
on average roughly \(D \times V_M\), and because its winding and its inertia smooth the
pulses, it turns as if fed that average voltage. The ESP32-S3 generates the PWM signals
with its LED PWM controller, at 20~kHz with 8-bit resolution, so the duty cycle is set
by a number from 0 to 255 (firmware files \texttt{pins\_esp32s3.h} and
\texttt{motor.cpp}). For forward motion the firmware applies PWM to \texttt{IN1} and holds
\texttt{IN2} low; for reverse it does the opposite, and it sets both inputs low to stop.
In the DRV8833's terms this is PWM with fast decay \cite{drv8833_datasheet}: during the
off part of each period the bridge is disabled and the motor coasts. Four PWM channels
are therefore needed for the two motors.

\paragraph{Why the DRV8833.} One chip drives both motors, which saves space and wiring on
the waffle plate. Its supply range, 2.7 to 10.8~V, covers the whole range of a 2S
lithium-ion pack, and its logic inputs accept the 3.3~V signals of the ESP32-S3
directly. It was also available locally at low cost.

The module we use is a small red breakout board that carries the chip and a 10~\textmu F
capacitor on the motor supply, and follows the pin layout of the common TB6612FNG
boards, with the three pins the DRV8833 does not need marked NC. The chip's marking is
not readable on our module, so we cannot tell from it whether it is the HTSSOP version,
whose exposed pad allows 1.5~A per bridge, or the plain TSSOP one, limited to 0.5~A
\cite{drv8833_datasheet}. We therefore take the lower figure, 0.5~A continuous per motor,
as the safe limit.

\paragraph{The driver we started with.} The robot was first wired with two DRV8871
modules, one per motor. With them, the motors never turned at all, while the same motors
ran with the DRV8833. We did not
measure the battery voltage at the time, so we cannot state the cause with certainty.
The DRV8871 datasheet points to a likely one: the chip is rated for a motor supply of
6.5 to 45~V, and its undervoltage lockout switches the bridge off when the supply falls
to 6.1~V typically (6.4~V at most), only restarting once it rises back to about
6.3--6.5~V \cite{drv8871_datasheet}. Our 2S pack is 7.4~V nominal
(2~\(\times\)~3.7~V, from the cell labels), only about 1~V above that threshold, and its
voltage falls as it discharges and drops further for a moment each time a motor starts
and draws a surge of current. A partly discharged pack could therefore keep the DRV8871
at or below its lockout threshold. The DRV8833, whose lockout is at 2.6~V, has no such
problem with 2S.

\subsection{Motors: TT gearmotors}
Each wheel is driven by a TT gearmotor, the yellow plastic motor found in most
hobby robot kits. It combines a small brushed DC motor with a plastic gearbox that
reduces the speed and multiplies the torque, and has an output shaft on both sides of
the gearbox. One side carries the wheel; the other is free, which is where the encoder
disc of Section~\ref{sec:encoders} is fitted.

Our motors carry no marking, and the shops we bought them from publish no
specifications. We therefore use the figures published by Adafruit for the same common
1:48 TT motor as a reference, given in Table~\ref{tab:ttmotor}; they describe the type
of motor, not measurements of our own units.

\begin{table}[htbp]
  \centering
  \caption{Reference figures for the 1:48 TT gearmotor (vendor bench measurements
           \cite{adafruit_tt_motor}, not measured on our motors).}
  \label{tab:ttmotor}
  \begin{tabular}{@{}lcc@{}}
    \toprule
    & at 3~V & at 6~V \\
    \midrule
    No-load speed (output shaft) & 120~rpm & 250~rpm \\
    No-load current & 150~mA & 160~mA \\
    Stall current & 1.1~A & 1.5~A \\
    Stall torque & 0.4~kg\,cm & 0.8~kg\,cm \\
    \bottomrule
  \end{tabular}
\end{table}

Three consequences follow for the design.

\paragraph{Voltage.} The motor is rated for 3 to 6~V \cite{adafruit_tt_motor}, while
the 2S pack gives 7.4~V nominally and about 8.4~V when fully charged (4.2~V per
cell, the usual full-charge voltage of a lithium-ion cell).
The firmware therefore never applies full duty: drive commands are scaled so that the
duty cycle does not exceed 204/255, or 80~\%. On a pack at 7.7~V, the value used in the
firmware comment, this gives about \(0.8 \times 7.7 \approx 6.2\)~V on average. When the
pack is fully charged the average rises to about 6.7~V, slightly above the rating; this
is a known limit of the current setting.

\paragraph{Current.} While running, a TT motor draws well under 0.5~A, the safe limit we
set for the driver in Section~\ref{sec:drv8833}, but at stall, or for an instant when it starts, it
can draw up to about 1.5~A at 6~V (Table~\ref{tab:ttmotor}). This is still below the
2~A peak the DRV8833 accepts and below its overcurrent trip point
\cite{drv8833_datasheet}. To soften the start-up surge, the firmware does not jump to
the requested speed but ramps the duty cycle by 5 steps (out of 255) every 10~ms
(\texttt{motor.cpp}), so going from rest to the 80~\% cap takes about 0.4~s.

\paragraph{Speed.} The wheels are standard 65~mm TT wheels (the diameter used in the
CAD model and in the firmware; test T1 in Chapter~3 checks it in practice). One turn of
the wheel moves the robot by \(\pi \times 65 \approx 204\)~mm. At the 250~rpm no-load
speed of Table~\ref{tab:ttmotor}, this is about 0.85~m/s, or 3.1~km/h. Under the weight
of the robot the real speed is lower, so the robot should not go much faster than about
3~km/h. This order of magnitude is a useful check on the speed the encoders report.

\paragraph{Why TT motors.} They are geared, so they turn slowly enough to drive a wheel
directly, cheap, easy to mount, and the dual shaft leaves room for an encoder. Above
all they were the motors available to us.

\subsection{Wheel speed sensors: LM393 slotted encoders}
\label{sec:encoders}
To report the speed of each wheel, the robot uses one slotted speed sensor per motor:
a small board with a U-shaped optical fork, an LM393 comparator and a digital output.

\paragraph{Working principle.} One arm of the fork holds an infrared LED and the other a
phototransistor facing it. A plastic disc with evenly spaced slots is mounted on the
free shaft of the TT motor so that its rim passes through the fork. As the wheel turns,
the disc alternately lets the infrared light through a slot and blocks it, and the
phototransistor's current rises and falls with it. This signal is slow and noisy at its
edges, so the board feeds it to an LM393 comparator, which compares it with a fixed
reference and gives a clean two-level output: one pulse per slot. The LM393's output
stage is an open-collector transistor that can only pull the line low, so the board
carries a pull-up resistor that defines the high level \cite{lm393_datasheet}. We power
the sensors from the ESP32-S3's 3.3~V supply, so that this high level is 3.3~V and the
output can go straight to a GPIO.

\paragraph{From pulses to speed.} With \(N\) slots on the disc and a wheel of diameter
\(D\), each pulse corresponds to a distance travelled by the wheel of
\begin{equation}
  \Delta s = \frac{\pi D}{N}.
  \label{eq:mm-per-pulse}
\end{equation}
Our discs have \(N = 20\) slots, counted on the real parts, and the wheels are 65~mm in
diameter, which gives \(\Delta s = \pi \times 65 / 20 \approx 10.2\)~mm per pulse
(\texttt{encoders.cpp}). The firmware first assumed 16 slots; this made every reported speed
25~\% too high (20/16) until the
slots were counted and the constant corrected. Each sensor output
is wired to its own GPIO (40 and 41) and triggers an interrupt on every rising edge.

\paragraph{First method: counting pulses.} The first version of the firmware counted the
pulses of each wheel and, every 100~ms, computed the speed as
\begin{equation}
  v = \frac{n \, \Delta s}{\Delta t},
\end{equation}
where \(n\) is the number of pulses counted during the interval \(\Delta t = 100\)~ms.
This turned out to be too coarse. At the robot's speeds a 100~ms window holds only a few
pulses, so one pulse more or less changes the result by
\(10.2~\text{mm} / 0.1~\text{s} \approx 102\)~mm/s, about 0.37~km/h, and the displayed speed
jumped between neighbouring values.

\paragraph{Current method: timing the pulses.} The firmware now measures the time between
pulses instead. Each interrupt also records the period \(T_i\) since the previous pulse,
taken from the microsecond timer. Every 100~ms, the speed is computed from the average
period of the \(n\) pulses received in that window:
\begin{equation}
  v_{\text{meas}} = \frac{n \, \Delta s}{\sum_{i=1}^{n} T_i}.
\end{equation}
Since the periods are timed to the microsecond, the result no longer moves in steps of
about 100~mm/s. Three details handle the edge cases. If no pulse arrives during a window, the
time elapsed since the last pulse gives an upper bound on the speed, so the reading falls
smoothly instead of dropping at once. If no pulse arrives for 300~ms, the wheel is
considered stopped and the speed is set to zero; this means the slowest speed the robot
can report is about \(10.2~\text{mm} / 0.3~\text{s} \approx 34\)~mm/s. And the first
pulse after a stop is not used, because its ``period'' is just the time the wheel was
standing still. Finally, the measured value is smoothed with a first-order filter,
\begin{equation}
  v_k = v_{k-1} + \alpha \,(v_{\text{meas}} - v_{k-1}), \qquad \alpha = 0.4,
\end{equation}
where a larger \(\alpha\) follows changes faster but lets more noise through.

The sensor has a single channel, so it sees that the wheel turns but not in which
direction. The firmware takes the sign from the command it is applying to that motor,
then sends both wheel speeds, in millimetres per second, to the controller every 100~ms.

\paragraph{Protection against false pulses.} Slotted sensors of this kind can produce
several rapid transitions as a slot edge passes, which would be counted as extra pulses.
The interrupt routine therefore ignores any pulse that arrives less than 1~ms after the
previous one on the same wheel. At the TT motor's no-load speed of 250~rpm
(Table~\ref{tab:ttmotor}), a 20-slot disc gives about 83 pulses per second, one every
12~ms, so this 1~ms window cannot hide a real pulse.

\paragraph{Why these sensors.} The board's fork fits around the rim of a disc on the TT
motor's second shaft, the sensors cost little, and they were what we had. They give
speed, not position or direction; a quadrature encoder would give both, at a higher cost.

\subsection{Ultrasonic sensor: HC-SR04}
\label{sec:hcsr04}
At the front of the robot, an HC-SR04 ultrasonic module measures the distance to
the nearest obstacle. It is mounted in a VBloc on the upper plate, facing forward.

\paragraph{Working principle.} The HC-SR04 measures distance by time of flight. When its
\texttt{TRIG} input receives a pulse of at least 10~\textmu s, its transmitter sends a
burst of eight ultrasonic cycles at 40~kHz. The sound travels to the obstacle, is
reflected, and comes back to the receiver. The module raises its \texttt{ECHO} output
when the burst leaves and lowers it when the echo returns, so the length \(t\) of the
\texttt{ECHO} pulse is the round-trip time \cite{hcsr04_datasheet}. Since the sound
covers the distance twice, the distance \(d\) to the obstacle is
\begin{equation}
  d = \frac{v \, t}{2},
  \label{eq:tof}
\end{equation}
where \(v\) is the speed of sound in air. The datasheet takes \(v = 340\)~m/s, which
gives the rule of thumb \(d\,[\text{cm}] = t\,[\text{\textmu s}] / 58\)
\cite{hcsr04_datasheet}. The value of \(v\) depends on the air temperature, so the
measured distance carries a small error when the air is much warmer or colder than
assumed. Table~\ref{tab:hcsr04} lists the module's main figures.

\begin{table}[htbp]
  \centering
  \caption{Main figures of the HC-SR04 module, from the manufacturer's datasheet
           \cite{hcsr04_datasheet}.}
  \label{tab:hcsr04}
  \begin{tabularx}{\textwidth}{@{}l X@{}}
    \toprule
    Parameter & Value \\
    \midrule
    Supply & 5~V, 15~mA \\
    Range & 2~cm to 4~m \\
    Measuring angle & 15\textdegree \\
    Trigger & pulse of at least 10~\textmu s \\
    Ultrasonic burst & 8 cycles at 40~kHz \\
    Recommended cycle & at least 60~ms between measurements \\
    \bottomrule
  \end{tabularx}
\end{table}

\paragraph{Reading it without blocking.} The maximum range sets how long an echo can
take: for an obstacle at 4~m the sound travels 8~m, which takes about
\(8 / 340 \approx 24\)~ms. A simple program would wait for the \texttt{ECHO} pulse to end,
but waiting up to 25~ms in the main loop would delay the motor ramp and the failsafe,
which run every 10~ms. The robot's firmware therefore times the echo with an interrupt
on GPIO~12: the rising edge stores the start time, the falling edge the pulse width,
and the loop never waits. Every 100~ms the firmware converts the last pulse width into a
distance with Equation~\eqref{eq:tof}, using \(v = 343\)~m/s (0.343~mm/\textmu s), sends
it to the controller in millimetres, and fires the next ping. A pulse longer than
25~ms, about 4.3~m, is reported as ``no object''. The 100~ms spacing also respects the
datasheet's advice to leave at least 60~ms between measurements, so that a late echo is
not mistaken for the next one.

\paragraph{Level shifting.} The HC-SR04 runs from 5~V, and its \texttt{ECHO} output
swings to that supply. The ESP32-S3, however, accepts at most its own supply plus
0.3~V on an input, about 3.6~V \cite{esp32s3_datasheet}, so \texttt{ECHO} cannot be
connected directly. A voltage divider of two resistors brings it down: \(R_1\) in series
between \texttt{ECHO} and GPIO~12, and \(R_2\) from GPIO~12 to ground. The pin then sees
\begin{equation}
  V_{\text{GPIO}} = V_{\text{ECHO}} \, \frac{R_2}{R_1 + R_2}.
\end{equation}
We use \(R_1 = 1\)~k\(\Omega\) and \(R_2 = 2\)~k\(\Omega\), so a 5~V echo gives 3.33~V,
within the input limit and well above the \(0.75 \times V_{DD}\), about 2.5~V, that the
ESP32-S3 needs to read a high level \cite{esp32s3_datasheet}. The \texttt{TRIG} input is
left without conversion. The datasheet only asks for a ``TTL pulse''
\cite{hcsr04_datasheet} and gives no threshold, but in practice the module triggers
reliably on the 3.3~V pulse from GPIO~11.

\paragraph{Why the HC-SR04.} It is a cheap, widely sold module that is easy to drive
with two pins, and it was available. Its limits are stated by the manufacturer: the beam
spreads over about 15\textdegree, and reliable readings need a target of at least
0.5~m\(^2\) with a fairly smooth surface \cite{hcsr04_datasheet}.

\subsection{Inertial sensor: MPU6050}
\label{sec:mpu6050}
The robot also carries an MPU6050 inertial measurement unit (IMU), which tells it how it
is tilted and how fast it turns. At this stage the data is sent to the controller and
displayed; it is the basis for later uses such as holding a straight line or detecting
that the robot has tipped over.

\paragraph{What it measures.} The MPU6050 from InvenSense combines a three-axis MEMS
gyroscope, which measures angular rate, and a three-axis MEMS accelerometer, which
measures acceleration including gravity, in one chip \cite{mpu6000_ps}. It runs from
2.375 to 3.46~V and is read over I\textsuperscript{2}C at up to 400~kHz, at the address
0x68 when its AD0 pin is low \cite{mpu6000_ps}. On the robot it shares an
I\textsuperscript{2}C bus, on GPIO~8 (SDA) and GPIO~9 (SCL), that is reserved for later
I\textsuperscript{2}C sensors as well \todo{module type (GY-521?) and which supply pin it is on}.

\paragraph{Configuration.} At start-up the firmware wakes the chip by clearing the
\texttt{SLEEP} bit of its power-management register, and selects the most sensitive
ranges: \(\pm 2\)~g for the accelerometer and \(\pm 250\)\textdegree/s for the gyroscope
\cite{mpu6000_rm}. In these ranges one unit of the 16-bit raw value is
\(1/16\,384\)~g and \(1/131\)\textdegree/s \cite{mpu6000_ps}, which are the scale factors
the firmware divides by. It then averages 200 gyroscope readings while the robot stands
still and stores them as offsets, because a real gyroscope does not read exactly zero at
rest; the robot must therefore not be moved for about a second after it is switched on.

\paragraph{From raw data to tilt.} Two estimates of the tilt are available, and each is
good where the other is weak. The accelerometer gives the roll and pitch angles directly
from the direction of gravity,
\begin{equation}
  \phi_{\text{acc}} = \operatorname{atan2}(a_y, a_z), \qquad
  \theta_{\text{acc}} = \operatorname{atan2}\!\left(-a_x, \sqrt{a_y^2 + a_z^2}\right),
\end{equation}
but this estimate is disturbed every time the robot accelerates or vibrates. The
gyroscope, integrated over time, follows fast changes cleanly but slowly drifts away
from the true angle. The firmware blends the two with a complementary filter, run on
every pass of the main loop:
\begin{equation}
  \phi_k = 0.96 \,(\phi_{k-1} + \omega_x \, \Delta t) + 0.04 \, \phi_{\text{acc}},
\end{equation}
and the same for the pitch \(\theta\) with \(\omega_y\). The gyroscope term dominates over
short times, while the small accelerometer share slowly pulls the estimate back and
cancels the drift. Every 50~ms (20~Hz) the robot sends the roll, the pitch and the three
angular rates to the controller, which shows them on its MPU6050 screen.

\paragraph{Why the MPU6050.} It is a common and cheap module that gives six axes over two
wires, and it was available to us. The yaw angle, the heading, cannot be corrected the
same way, because gravity gives no reference for it; it would need a magnetometer, and
the firmware does not compute it.

\subsection{Power supply}
\label{sec:robot-power}
\paragraph{The battery.} The robot runs from two 18650 lithium-ion cells, the
cylindrical 18~mm by 65~mm format used in laptop packs and power banks. Ours are METAMA
cells labelled 3.7~V and 2500~mAh; we found no datasheet for them, so the label values
are all we use. The two cells sit in a two-slot holder and are connected in series, the
positive terminal of one to the negative terminal of the other, which makes a 2S pack:
\begin{itemize}
  \item nominal voltage \(2 \times 3.7 = 7.4\)~V, rising to about 8.4~V when fully
        charged (4.2~V per cell, the usual full-charge voltage of a lithium-ion cell);
  \item capacity 2500~mAh, since cells in series share the same current;
  \item nominal energy about \(7.4~\text{V} \times 2.5~\text{Ah} \approx 18.5\)~Wh.
\end{itemize}
How long the robot runs on one charge is measured in Chapter~3 (test T2) rather than
estimated here. The cells are removed from the robot and charged in a separate two-slot
lithium-ion charger.

\paragraph{Distribution.} The pack feeds two branches, as shown in
Figure~\ref{fig:architecture}. The first goes straight to the motor supply pin
\(V_M\) of the DRV8833, whose 2.7 to 10.8~V range covers the whole 2S range
(Section~\ref{sec:drv8833}). The second goes to a step-down converter that produces the
5~V for the ESP32-S3 board, whose own regulator then makes the 3.3~V the module and the
LM393 sensors use. The HC-SR04, which needs 5~V, takes it from the board's 5~V pin, that
is from the converter's output (\texttt{pins\_esp32s3.h}).

\paragraph{The step-down converter.} A pack that reaches 8.4~V cannot be connected to the
board's 5~V pin directly. A linear regulator would turn the difference into heat; a
switching \emph{buck} converter instead switches the input on and off at high frequency
and smooths the result with an inductor and a capacitor, so it wastes far less. We use a
Mini360 module, a very small buck board built around the MP2307 controller, whose output
voltage is set with an on-board trimmer. Before the ESP32-S3 was connected, we fed
the module 7.4~V from a bench power supply and turned the trimmer until a multimeter
read 5~V at its output. Its main figures, from a module description \cite{mini360_c101}, are an input
of 4.75 to 23~V, an output of 1 to 17~V, and 1.8~A continuous (3~A peak). The ESP32-S3
draws at most about 340~mA, while transmitting over Wi-Fi \cite{esp32s3_datasheet}, so
the converter works far below its limit.

\paragraph{What is missing.} The pack has no battery management system (BMS) and no
fuse, and the main switch is planned but not fitted yet. A BMS would cut the load when a
cell is over-discharged and stop overcharging, and a fuse would protect the wiring from
a short circuit; without them, these protections rely on the user, who must not run the
pack flat and must charge the cells in the dedicated charger. One of our cells failed
during development and was replaced; we did not identify the cause. Adding a 2S BMS, a
fuse and the switch is the first item of future work on the hardware.

Decoupling is also minimal. TI asks for at least 10~\textmu F close to the DRV8833's
\(V_M\) pin, which the module provides, and recommends adding a local bulk capacitor,
rated above the supply voltage, whose size is to be found by testing
\cite{drv8833_datasheet}. No bulk capacitor has been added yet, nor the small ceramic
capacitors often soldered across a DC motor's terminals to absorb the noise of its
brushes. Both are cheap to add: a 470~\textmu F, 16~V electrolytic capacitor across the
driver's \(V_M\) and GND pins, and a 0.1~\textmu F capacitor on each motor.

\section{The controller}
\label{sec:controller}

The controller is the handheld unit the user drives the robot with. It reads a joystick
and four buttons, shows menus and the robot's data on a small screen, and talks to the
robot over ESP-NOW. Its parts are shown in the upper half of
Figure~\ref{fig:architecture}, and its pin assignment in Table~\ref{tab:con-pins}.

\subsection{Microcontroller: ESP32 DevKit}
The controller uses a classic ESP32, the predecessor of the chip on the robot. Its main
figures are a dual-core Xtensa LX6 processor at up to 240~MHz, 520~KB of SRAM, 2.4~GHz
Wi-Fi (IEEE 802.11b/g/n), Bluetooth~4.2, two 12-bit analog-to-digital converters and 34
GPIOs, six of which are inputs only \cite{esp32_datasheet}. The board is a generic
30-pin development board of the ``ESP32 DevKit V1'' type, with a USB-C connector and a
CP2102 USB-to-serial converter for programming. Its module carries no marking, so we
cannot name it exactly; like the robot's board, it takes 5~V on its \texttt{VIN} pin and
regulates it to 3.3~V.

\paragraph{Why an ESP32.} The controller's job is light: read a few inputs, draw a
128 \(\times\) 64 screen and exchange small packets with the robot. A classic ESP32 does
this easily and it is the board we already had. Since it uses the same radio and the same
ESP-NOW protocol as the ESP32-S3, the two ends speak to each other without any extra
hardware. A smaller and cheaper chip, such as the single-core ESP32-C3, would probably
be enough and is left as a possible improvement.

\begin{table}[htbp]
  \centering
  \caption{GPIO assignment on the controller (firmware file \texttt{Bloc-Con/include/config.h}).}
  \label{tab:con-pins}
  \begin{tabularx}{\textwidth}{@{}l l X@{}}
    \toprule
    GPIO & Connected to & Note \\
    \midrule
    34, 35 & joystick X, Y & ADC1, input-only pins \\
    12, 13, 14, 27 & buttons UP, DOWN, FORWARD, BACK & to ground, internal pull-ups;
                     12 is a strapping pin (see text) \\
    21, 22 & OLED SDA, SCL & I\textsuperscript{2}C \\
    \bottomrule
  \end{tabularx}
\end{table}

\paragraph{A pin to handle with care.} GPIO~12 is one of the ESP32's strapping pins: its
level at reset, through the bit called MTDI, selects the voltage of the flash memory. Its
default, set by an internal pull-down, is 0, which keeps the 3.3~V flash supply
\cite{esp32_datasheet}. The UP button is connected to it, which is safe because the button
can only pull the pin to ground and the firmware enables the internal pull-up only after
boot. An external pull-up on this pin, however, could select the wrong flash voltage and
stop the board from starting.

\subsection{Human interface: joystick, buttons and display}

\paragraph{Joystick.} The KY-023 module is a two-axis thumb joystick: each axis moves the
wiper of a potentiometer connected between the supply and ground, so the wiper voltage
varies with the position of the stick and rests near mid-supply when the stick is
released. The module also has a push switch, which the controller does not use. Although
its supply pin is labelled +5~V, we power it from 3.3~V so that its outputs never exceed
what the ESP32's inputs accept. The two wipers go to GPIO~34 and~35, which belong to ADC1.
This matters because the radio is on all the time and ADC2 is shared with the Wi-Fi
driver \cite{espidf_adc}, while ADC1 is not. The ADC converts each voltage into a 12-bit
number, from 0 to 4095. Its accuracy is limited: the datasheet gives a non-linearity of up
to 12~LSB and advises averaging several samples \cite{esp32_datasheet}, which is what the
firmware does (Chapter~3).

\paragraph{Buttons.} Four push buttons, UP, DOWN, FORWARD (select) and BACK, navigate the
menus. Each connects its GPIO to ground when pressed. The pins use the ESP32's internal
pull-up resistors, so no external resistor is needed: a released button reads high, a
pressed one reads low.

\paragraph{Display.} The screen is a 0.96-inch monochrome OLED of 128 \(\times\) 64 pixels,
driven by a Solomon Systech SSD1306 controller. It is connected over
I\textsuperscript{2}C on GPIO~21 (SDA) and GPIO~22 (SCL), at the address 0x3C, the one the
SSD1306 answers to when its address-select pin is low \cite{ssd1306_datasheet}, and is
powered from 3.3~V, within the controller's 1.65 to 3.3~V logic supply range
\cite{ssd1306_datasheet}. The firmware draws each screen in a frame buffer in the ESP32's
memory, \(128 \times 64 / 8 = 1024\) bytes, and sends it to the display in one transfer.
OLED pixels emit their own light, so the screen needs no backlight. During development we found the small 0.96-inch screens hard to read at arm's
length, so a larger 2.42-inch display of the same resolution, driven by the closely
related SSD1309, is planned. The firmware already supports it through a single
compile-time switch (\texttt{USE\_SSD1309}), and the wiring is the same.

\subsection{Power supply}
\label{sec:con-power}
The controller is powered by a single 18650 cell, a QHTC cell labelled 2000~mAh and
4.2~V. It is a cheaper cell than the robot's, which is acceptable because the controller
draws far less current than two motors. The power chain has three stages, shown in
Figure~\ref{fig:architecture}.

\paragraph{Charging and protection: TP4056 module.} The cell stays in the controller and
is charged through a TP4056 module with a USB-C input. The TP4056 is a linear charger for
single lithium-ion cells. It charges at constant current, set by a programming resistor,
then at constant voltage once the cell reaches 4.2~V (within 1.5~\%), and charges a deeply
discharged cell gently until it rises above 2.9~V \cite{tp4056_datasheet}. Our module
carries no reference; it is identical in layout to a USB-C TP4056 module sold locally,
whose listing we use for its figures \cite{esp32lab_tp4056}: a maximum charge current of
1000~mA.

Unlike the robot's pack, the controller's cell is protected. The module has separate
terminals for the battery (B+, B$-$) and for the load (OUT+, OUT$-$), and according to
that listing it disconnects the load when the cell falls to 2.5~V or when the current
exceeds 3~A \cite{esp32lab_tp4056}. The power switch is placed on the load side, after
the charger, so the cell can be charged with the controller switched off.

\paragraph{Boost: MT3608.} A lithium-ion cell stays below 4.2~V even when fully charged,
under the 5~V the ESP32 board expects on its \texttt{VIN} pin. A step-up (boost) converter raises
it. The MT3608 accepts 2 to 24~V at its input, switches at 1.2~MHz and regulates its
output so that its feedback pin sits at 0.6~V, which gives
\begin{equation}
  V_{\text{OUT}} = 0.6~\text{V} \times \left(1 + \frac{R_1}{R_2}\right),
\end{equation}
where \(R_1\) and \(R_2\) form the divider from the output to the feedback pin
\cite{mt3608_datasheet}. On the module, \(R_1\) is a trimmer. The datasheet allows outputs
up to 28~V, so the trimmer must be set to 5~V with a multimeter before the ESP32 is connected.

\paragraph{Load.} The ESP32 draws up to about 240~mA while transmitting over Wi-Fi at
full power \cite{esp32_datasheet}, and the screen and joystick add little to this. A
boost converter draws more current from the cell than it delivers, by the ratio of the
voltages and its efficiency, so the controller's runtime on one charge is measured
rather than estimated (Chapter~3).

\section{Wireless link: ESP-NOW}
\label{sec:espnow}
The controller and the robot communicate over ESP-NOW. This section explains what the
protocol is, why it was chosen, and what its use in BLOC-BOT does and does not guarantee.
The exact packet formats are given with the firmware in Chapter~3.

\paragraph{What ESP-NOW is.} ESP-NOW is, in Espressif's words, ``a kind of
connectionless Wi-Fi communication protocol that is defined by Espressif''
\cite{espnow_idf}. It uses the same 2.4~GHz radio as Wi-Fi, but instead of joining a
network, a device simply sends short frames addressed to another device's MAC address,
or to the broadcast address FF:FF:FF:FF:FF:FF so that every listening device receives
them. No access point and no connection set-up are needed. A frame carries at most
250~bytes of data in version~1.0 of the protocol (1470~bytes in version~2.0), and a device
can register up to 20 peers \cite{espnow_idf}. On reception, the firmware is notified by a
callback function that receives the sender's address and the data.

\paragraph{Why ESP-NOW and not Wi-Fi or Bluetooth.} The link has a simple job: carry a few
bytes of commands one way and a few bytes of measurements the other, several times per
second, between two boards that are both Espressif chips.
\begin{itemize}
  \item \textbf{Wi-Fi} in the usual sense would require one board to act as an access
        point and the other to join it before any data could flow, and then an
        IP-based protocol on top. This is more set-up than two boards exchanging small
        packets need.
  \item \textbf{Bluetooth}: the two chips do not even share the same versions. The
        controller's ESP32 supports Bluetooth~4.2 in both its Classic and Low Energy
        forms \cite{esp32_datasheet}, but the robot's ESP32-S3 supports only Bluetooth~5
        Low Energy \cite{esp32s3_datasheet}. A Bluetooth LE link would work, but it needs
        a connection between a central and a peripheral device and a description of the
        exchanged data as services and characteristics.
  \item \textbf{ESP-NOW} needs neither an access point nor a connection: each board can
        send as soon as the radio is started, and the same few lines of code run on both
        chips. The 250-byte limit is far above what BLOC-BOT needs, since its largest
        packet is 11~bytes.
\end{itemize}
The choice was also practical: the protocol was suggested to us during development, and
it suits a system in which every board is an ESP32.

\paragraph{How BLOC-BOT uses it.} Both boards send to the broadcast address, without
pairing and without encryption, on the radio's current channel. Every packet starts with a
one-byte identifier, and the receiver accepts a packet only if both this identifier and
the packet's length match what it expects; anything else is ignored. The controller sends
drive commands every 50~ms while its Drive screen is open, and the robot sends its wheel
speeds and the measured distance every 100~ms and its tilt every 50~ms.

\paragraph{What the link does not guarantee.} Three limits follow from these choices. The
first is handled by the design; the other two remain open.
\begin{itemize}
  \item \textbf{Delivery.} A successful send only means that the frame was received at the
        radio's MAC layer; Espressif states that reception by the application is not
        guaranteed \cite{espnow_idf}. BLOC-BOT therefore does not rely on any single
        packet: commands are repeated every 50~ms, and the robot stops its motors if none
        arrives for 300~ms, that is after about six missed commands in a row (requirement
        FR3).
  \item \textbf{Security and interference.} Broadcast frames cannot be encrypted
        \cite{espnow_idf}, and any ESP-NOW device in range that sends a packet with the
        same identifier and length would be obeyed. Two BLOC-BOTs used side by side would
        therefore interfere. Pairing each controller with its robot by MAC address, and
        enabling ESP-NOW's encryption, which is available for paired devices
        \cite{espnow_idf}, would remove this limit and is left as future work.
  \item \textbf{Range.} We have not measured the range of the link, and the documentation
        we used gives no figure for it, so the thesis states none.
\end{itemize}

\section{The modular platform}
\label{sec:platform}
\plan{Written last, short. Waffle plate (210 x 210 x 10~mm PETG, hole pattern
3.2/4.5/6.0~mm), blocs HBloc / VBloc / HBloc\_L, pillars. Old plate (printed) vs new plate
(designed): print time and mass. One or two figures, the print cost table.}

\section{Conclusion}
\plan{Short bridge to Chapter 3 (realisation and tests).}
